% GENERATED FILE. Edit canonical lessons, then run npm run generate:books.
% canonical-source-hash: fnv1a64:3da99aedbe7ab39c
% canonical-lessons: MW-C242-jalno, MW-C242-bujhno, MW-C242-phutno, MW-C242-thikro, MW-C242-dubno

\chapter{Fire, Water and a Broken Pot}
\label{ch:mw-fire-water-and-a-broken-pot}
\hlchaptermodality{242}

\begin{chapteropening}
\textbf{By the end of this chapter:} \emph{I can say what got burnt, went out, broke or sank.}

Complete the last lesson of chapter 242: I can say what got burnt, went out, broke or sank.
\end{chapteropening}

\section[\texorpdfstring{\mw{जळणो} (jaḷṇo)}{जळणो (jaḷṇo)}]{\mw{जळणो} (jaḷṇo) — to get burnt; to burn}

\label{lesson:MW-C242-jalno}



\begin{quote}
Before the new one: say the Marwadi for to get broken, then the Marwadi for to slip.
\end{quote}

\subsection*{You'll want to know: \mw{जळणो}}
\textbf{\mw{जळणो}} — \emph{jaḷṇo} — \textquotedblleft{}to get burnt; to burn\textquotedblright{}.

\begin{quote}
\mw{रसोई} \mw{मांय} \mw{म्हारो} \mw{हाथ} \mw{जळ्यो।}
\end{quote}

\begin{quote}
\emph{rasoī mā̃y mhāro hāth jaḷyo.}
\end{quote}

\begin{quote}
My hand got burnt in the kitchen.
\end{quote}

Past (it …): \mw{वो} \mw{पापड़} \mw{जळ्यो}, \emph{vo pāpaṛ jaḷyo}.

\begin{cousinweb}
\textbf{\mw{जळणो}} (\emph{jaḷṇo}) is inherited from Sanskrit \emph{jvalati}, \textquotedblleft{}burns brightly\textquotedblright{}; CDIAL lists Old Marwari \emph{jalaï}, \textquotedblleft{}burns\textquotedblright{}, under it (CDIAL 5306).
\end{cousinweb}

\begin{grammarlens}[title={What you've built}]
One more word for saying what happened.
\end{grammarlens}

\subsection*{Your turn}
\begin{itemize}
\raggedright
  \item \emph{Say it:} \emph{jaḷṇo}
  \item \emph{Say it:} \emph{jaḷṇo}, once more
  \item \emph{Say it:} \emph{rasoī mā̃y mhāro hāth jaḷyo.} — My hand got burnt in the kitchen.
  \item \emph{Recall:} say the Marwadi for to get broken, then the Marwadi for to slip, then say \emph{jaḷṇo} again
\end{itemize}

\begin{tcolorbox}[breakable,colback=teal!4,colframe=teal!35!black,title={Before you move on}]
What does \mw{जळणो} mean? (\textquotedblleft{}To get burnt; to burn\textquotedblright{}.) Say it once more.
\end{tcolorbox}
\section[\texorpdfstring{\mw{बुझणो} (bujhṇo)}{बुझणो (bujhṇo)}]{\mw{बुझणो} (bujhṇo) — to go out (of a flame)}

\label{lesson:MW-C242-bujhno}



\begin{quote}
Before the new one: say the Marwadi for to slip, then the Marwadi for to get burnt.
\end{quote}

\subsection*{You'll want to know: \mw{बुझणो}}
\textbf{\mw{बुझणो}} — \emph{bujhṇo} — \textquotedblleft{}to go out (of a flame)\textquotedblright{}.

\begin{quote}
\mw{रात} \mw{मांय} \mw{दीयो} \mw{बुझ्यो।}
\end{quote}

\begin{quote}
\emph{rāt mā̃y dīyo bujhyo.}
\end{quote}

\begin{quote}
The lamp went out in the night.
\end{quote}

Past (it …): \mw{वो} \mw{चूल्हो} \mw{बुझ्यो}, \emph{vo chūlho bujhyo}.

\begin{cousinweb}
\textbf{\mw{बुझणो}} (\emph{bujhṇo}) is inherited from a reconstructed \emph{*vijjhāyati}, \textquotedblleft{}is burnt out, is extinguished\textquotedblright{}; CDIAL lists the Old Marwari partner \emph{bujhāvaï}, \textquotedblleft{}to extinguish\textquotedblright{}, under it (CDIAL 11703).
\end{cousinweb}

\begin{grammarlens}[title={What you've built}]
One more word for saying what happened.
\end{grammarlens}

\subsection*{Your turn}
\begin{itemize}
\raggedright
  \item \emph{Say it:} \emph{bujhṇo}
  \item \emph{Say it:} \emph{bujhṇo}, once more
  \item \emph{Say it:} \emph{rāt mā̃y dīyo bujhyo.} — The lamp went out in the night.
  \item \emph{Recall:} say the Marwadi for to slip, then the Marwadi for to get burnt, then say \emph{bujhṇo} again
\end{itemize}

\begin{tcolorbox}[breakable,colback=teal!4,colframe=teal!35!black,title={Before you move on}]
What does \mw{बुझणो} mean? (\textquotedblleft{}To go out (of a flame)\textquotedblright{}.) Say it once more.
\end{tcolorbox}
\section[\texorpdfstring{\mw{फूटणो} (phūṭṇo)}{फूटणो (phūṭṇo)}]{\mw{फूटणो} (phūṭṇo) — to crack open; to burst}

\label{lesson:MW-C242-phutno}



\begin{quote}
Before the new one: say the Marwadi for to get burnt, then the Marwadi for to go out (of a flame).
\end{quote}

\subsection*{You'll want to know: \mw{फूटणो}}
\textbf{\mw{फूटणो}} — \emph{phūṭṇo} — \textquotedblleft{}to crack open; to burst\textquotedblright{}.

\begin{quote}
\mw{म्हारो} \mw{मटको} \mw{फूट्यो।}
\end{quote}

\begin{quote}
\emph{mhāro maṭko phūṭyo.}
\end{quote}

\begin{quote}
My big clay pot cracked.
\end{quote}

Past (it …): \mw{वो} \mw{नारेळ} \mw{फूट्यो}, \emph{vo nāreḷ phūṭyo}.

\begin{cousinweb}
\textbf{\mw{फूटणो}} (\emph{phūṭṇo}) is inherited from a reconstructed \emph{*sphuṭyati}, \textquotedblleft{}is burst open\textquotedblright{}. CDIAL gives no Marwari form in this entry; it lists Hindi \emph{phūṭnā}, \textquotedblleft{}to split, burst\textquotedblright{} (CDIAL 13845).
\end{cousinweb}

\begin{grammarlens}[title={What you've built}]
One more word for saying what happened.
\end{grammarlens}

\subsection*{Your turn}
\begin{itemize}
\raggedright
  \item \emph{Say it:} \emph{phūṭṇo}
  \item \emph{Say it:} \emph{phūṭṇo}, once more
  \item \emph{Say it:} \emph{mhāro maṭko phūṭyo.} — My big clay pot cracked.
  \item \emph{Recall:} say the Marwadi for to get burnt, then the Marwadi for to go out (of a flame), then say \emph{phūṭṇo} again
\end{itemize}

\begin{tcolorbox}[breakable,colback=teal!4,colframe=teal!35!black,title={Before you move on}]
What does \mw{फूटणो} mean? (\textquotedblleft{}To crack open; to burst\textquotedblright{}.) Say it once more.
\end{tcolorbox}
\section[\texorpdfstring{\mw{ठीकरो} (ṭhīkro)}{ठीकरो (ṭhīkro)}]{\mw{ठीकरो} (ṭhīkro) — a potsherd; a shard of a broken pot}

\label{lesson:MW-C242-thikro}



\begin{quote}
Before the new one: say the Marwadi for to go out (of a flame), then the Marwadi for to crack open.
\end{quote}

\subsection*{You'll want to know: \mw{ठीकरो}}
\textbf{\mw{ठीकरो}} — \emph{ṭhīkro} — \textquotedblleft{}a potsherd; a shard of a broken pot\textquotedblright{}.

\begin{quote}
\mw{ठीकरो} \mw{तीखो} \mw{है।}
\end{quote}

\begin{quote}
\emph{ṭhīkro tīkho hai.}
\end{quote}

\begin{quote}
The shard of the pot is sharp.
\end{quote}

A masculine noun.

\begin{cousinweb}
\textbf{\mw{ठीकरो}} (\emph{ṭhīkro}) is inherited from a reconstructed \emph{*ṭhīkkara}, \textquotedblleft{}potsherd\textquotedblright{}; CDIAL lists Marwari \emph{ṭhīkro}, \textquotedblleft{}potsherd\textquotedblright{}, under it (CDIAL 5504).
\end{cousinweb}

\begin{grammarlens}[title={What you've built}]
One more word for saying what happened.
\end{grammarlens}

\subsection*{Your turn}
\begin{itemize}
\raggedright
  \item \emph{Say it:} \emph{ṭhīkro}
  \item \emph{Say it:} \emph{ṭhīkro}, once more
  \item \emph{Say it:} \emph{ṭhīkro tīkho hai.} — The shard of the pot is sharp.
  \item \emph{Recall:} say the Marwadi for to go out (of a flame), then the Marwadi for to crack open, then say \emph{ṭhīkro} again
\end{itemize}

\begin{tcolorbox}[breakable,colback=teal!4,colframe=teal!35!black,title={Before you move on}]
What does \mw{ठीकरो} mean? (\textquotedblleft{}A potsherd; a shard of a broken pot\textquotedblright{}.) Say it once more.
\end{tcolorbox}
\section[\texorpdfstring{\mw{डूबणो} (ḍūbṇo)}{डूबणो (ḍūbṇo)}]{\mw{डूबणो} (ḍūbṇo) — to sink; to drown}

\label{lesson:MW-C242-dubno}



\begin{quote}
Before the new one: say the Marwadi for to crack open, then the Marwadi for a potsherd.
\end{quote}

\subsection*{You'll want to know: \mw{डूबणो}}
\textbf{\mw{डूबणो}} — \emph{ḍūbṇo} — \textquotedblleft{}to sink; to drown\textquotedblright{}.

\begin{quote}
\mw{नदी} \mw{मांय} \mw{नाव} \mw{डूबी।}
\end{quote}

\begin{quote}
\emph{nadī mā̃y nāv ḍūbī.}
\end{quote}

\begin{quote}
The boat sank in the river.
\end{quote}

Past (it …): \mw{वो} \mw{भाटो} \mw{डूब्यो}, \emph{vo bhāṭo ḍūbyo}.

\begin{cousinweb}
\textbf{\mw{डूबणो}} (\emph{ḍūbṇo}) is inherited from a reconstructed \emph{*ḍubb}, \textquotedblleft{}sink\textquotedblright{}. CDIAL gives no Marwari form in this entry; it lists Hindi \emph{ḍūbnā}, \textquotedblleft{}to sink\textquotedblright{} (CDIAL 5561).
\end{cousinweb}

\begin{grammarlens}[title={What you've built}]
One more word for saying what happened.
\end{grammarlens}

\subsection*{Your turn}
\begin{itemize}
\raggedright
  \item \emph{Say it:} \emph{ḍūbṇo}
  \item \emph{Say it:} \emph{ḍūbṇo}, once more
  \item \emph{Say it:} \emph{nadī mā̃y nāv ḍūbī.} — The boat sank in the river.
  \item \emph{Recall:} say the Marwadi for to crack open, then the Marwadi for a potsherd, then say \emph{ḍūbṇo} again
\end{itemize}

\begin{tcolorbox}[breakable,colback=teal!4,colframe=teal!35!black,title={Before you move on}]
What does \mw{डूबणो} mean? (\textquotedblleft{}To sink; to drown\textquotedblright{}.) Say it once more.
\end{tcolorbox}
